\VerdierPhysicalPageAnchor{0038}%
\VerdierSemanticAnchor{ch1-1-2-4-isomorphisms-cont}
and isomorphisms

\[
  m(g)\colon F_eT(g)\longrightarrow T'(g)F_e,
  \qquad g\in G,
\]
such that the diagrams (1.2.3.2) commute.

\VerdierSemanticAnchor{ch1-1-2-4-characterization}
A $G^\circ$-functor is thus a functor “commuting up to isomorphism
with the actions of $G^\circ$”. We shall use the notation $(F_e,m)$
to denote it.

\VerdierSemanticAnchor{ch1-1-2-4-strict}
When for every $g\in G$ the isomorphisms $m(g)$ are identity
isomorphisms, which in particular implies that for every $g\in G$ one
has the equalities $F_eT(g)=T'(g)F_e$, the $G^\circ$-functor is said to
be \emph{strict}.

\VerdierSemanticAnchor{ch1-1-2-5}
\noindent\textbf{1.2.5.}\quad
Let $\mathcal C$, $\mathcal C'$, and $\mathcal C''$ be three graded
categories, and let
\[
  \mathcal C\xrightarrow{F}\mathcal C'
  \xrightarrow{F'}\mathcal C'',\qquad F''=F'F,
\]
be three graded functors. Choose normalized cleavages on $\mathcal C$,
$\mathcal C'$, and $\mathcal C''$. From the preceding discussion one
obtains three $G^\circ$-categories
$(\mathcal C_e,G^\circ,T,c)$,
$(\mathcal C'_e,G^\circ,T',c')$, and
$(\mathcal C''_e,G^\circ,T'',c'')$, and three $G^\circ$-functors
$(F_e,m)$, $(F'_e,m')$, and $(F''_e,m'')$. One immediately checks the
relations

\VerdierSemanticAnchor{ch1-1-2-5-eq1}
\begin{equation}
  \begin{cases}
    F''_e=F'_eF_e,\\[0.5em]
    m''(g)=\bigl(m'(g)\star F_e\bigr)
      \bigl(F'_e\star m(g)\bigr),\qquad g\in G.
  \end{cases}
  \tag{1.2.5.1}
\end{equation}

\VerdierSemanticAnchor{ch1-1-2-6}
\noindent\textbf{Definition 1.2.6.}\quad
Let $(\mathcal C_e,G^\circ,T,c)$,
$(\mathcal C'_e,G^\circ,T',c')$, and
$(\mathcal C''_e,G^\circ,T'',c'')$ be three $G^\circ$-categories and
let

\[
\begin{tikzcd}[column sep=large]
(\mathcal C_e,G^\circ,T,c)
  \arrow[r,"{(F_e,m)}"]
&
(\mathcal C'_e,G^\circ,T',c')
  \arrow[r,"{(F_e^{\prime},m^{\prime})}"]
&
(\mathcal C''_e,G^\circ,T'',c'')
\end{tikzcd}
\]
be two $G^\circ$-functors.

\VerdierSemanticAnchor{ch1-1-2-6-composite}
The equalities (1.2.5.1) define a $G^\circ$-functor called the
\emph{composite functor} of the $G^\circ$-functors $(F_e,m)$ and
$(F'_e,m')$.

\VerdierSemanticAnchor{ch1-1-2-7}
\noindent\textbf{1.2.7.}\quad
Let $F$ and $F'$ be two graded functors between two $G$-graded
categories $\mathcal C$ and $\mathcal C'$, and let
$u\colon F\to F'$ be a morphism of functors of degree $e$. Choose
normalized cleavages on $\mathcal C$ and $\mathcal C'$. One obtains
two $G^\circ$-categories $(\mathcal C_e,G^\circ,T,c)$ and
$(\mathcal C'_e,G^\circ,T',c')$, and two $G^\circ$-functors
$(F_e,m)$ and $(F'_e,m')$. The morphism $u$ then determines a morphism
of functors

\VerdierSemanticAnchor{ch1-1-2-7-eq1}
\begin{equation}
  u_e\colon F_e\longrightarrow F'_e
  \tag{1.2.7.1}
\end{equation}
such that the following diagram commutes:

\VerdierSemanticAnchor{ch1-1-2-7-eq2}
\begin{equation}
\begin{tikzcd}[column sep=huge,row sep=large]
F_eT(g)
  \arrow[r,"u_e\star T(g)"]
  \arrow[d,"m(g)",swap]
&
F_e^{\prime}T(g)
  \arrow[d,"m^{\prime}(g)"]
\\
T^{\prime}(g)F_e
  \arrow[r,"T^{\prime}(g)\star u_e",swap]
&
T^{\prime}(g)F_e^{\prime}\mathpunct{.}
\end{tikzcd}
\tag{1.2.7.2}
\end{equation}
